\begin{figure}[H]
\centering
\begin{tikzpicture}[font=\small,>=Latex]
  \begin{scope}
    \draw[->,gray] (-.2,0)--(4.1,0) node[right,black] {\(X\)};
    \fill (0,0) circle (1.5pt) node[below=4pt] {\(0\)};
    \foreach \position/\radius/\disklabel in {3.2/.20/1,1.6/.10/2,.8/.05/3} {
      \draw[fill=orange!18] (\position,0) circle (\radius);
      \node[above=10pt] at (\position,0) {\(\mathcal D_{\disklabel}\)};
    }
    \fill (.4,0) circle (.025);
    \fill (.2,0) circle (.0125);
    \draw[->,blue!65!black] (2.6,-.50)--(.6,-.50);
    \node at (2,-1.05) {(a) Disks approaching the origin};
  \end{scope}
  \begin{scope}[shift={(7.1,.18)}]
    \draw[fill=orange!18] (0,0) circle (.67);
    \draw[gray,dashed] (0,0) circle (.90);
    \draw[gray,dashed] (0,0) circle (1.23);
    \foreach \angle/\innerrotation/\outerrotation in {
      10/0/60,50/30/0,90/60/30,130/15/75,170/45/15,
      210/75/45,250/0/30,290/60/0,330/30/60} {
      \begin{scope}[shift={(\angle:.923)},rotate=\innerrotation]
        \draw[blue!65!black,fill=blue!8] (-.033,-.033) rectangle (.033,.033);
      \end{scope}
      \begin{scope}[shift={(\angle:1.25)},rotate=\outerrotation]
        \draw[blue!65!black,fill=blue!8] (-.07,-.07) rectangle (.07,.07);
      \end{scope}
    }
    \node at (0,0) {\(\mathcal D_j\)};
    \node at (0,-1.55) {(b) Patches approaching a disk};
  \end{scope}
\end{tikzpicture}
\caption{The two accumulation levels, schematically and not to scale.
The disks in (a) are disjoint and avoid the origin. In (b), disjoint
full patches occupy successive annuli outside one disk. Only a few
patches and annuli are shown. Patch orientations are independent of
their positions around the disk; the actual construction approaches
every boundary point in every direction of a fixed finite rotation net.}
\label{fig:accumulation}
\end{figure}
